%% kapa: exam=1 round=22 year=2011 period=1 form=A subject=civil_law subject_ko=민법 seq=27 number=27
%% source: Q-Net 감정평가사 기출문제 (cst00309) artlSeq=5036790 file=extracted/1008/5036790/2011년 제22회 감정평가사 1차시험 기출문제/1교시.hwp
%% conversion: kordoc 4.24.51 -> tools/gen1.py (auto); review=proofread-sonnet; flags=none
\begin{question}{27}
법인의 이사에 관한 설명으로 옳지 않은 것은? (다툼이 있으면 판례에 의함)
\begin{choices}
  \item 이사의 임면에 관한 규정은 정관의 필요적 기재사항이고, 이사의 성명과 주소는 등기사항이다.
  \item 이사는 정관 또는 총회의 결의로 금지하지 않은 사항에 한하여 타인으로 하여금 특정한 행위를 대리하게 할 수 있다.
  \item 이사는 정관에 특별한 제한이 없는 한 자유롭게 사임할 수 있으며, 사임행위는 상대방 없는 단독행위이다.
  \item 법인에는 반드시 이사를 두어야 하며, 이사의 수는 정관에서 임의로 정할 수 있다.
  \item 이사의 행위가 사리(私利)를 도모하기 위한 것이더라도 외관상 객관적으로 직무에 관한 행위라고 인정할 수 있는 것이라면 ‘법인의 불법행위책임’(제35조)에서의 직무에 관한 행위에 해당한다.
\end{choices}
\end{question}
\answer{3}
